\documentclass[11pt]{article}
\usepackage{varnernotes}

\newcommand{\R}{\mathbb{R}}

\begin{document}

\lecturenotes{8b}{SIM-Based Risky and Risk-Free Portfolio Allocation}{Fall 2026}{October 15, 2026}{Jeffrey D. Varner}

\learningobjectives
  {Construct SIM portfolio inputs}
  {Build expected-return and positive-semidefinite covariance estimates from market, beta, alpha, and residual-risk estimates.}
  {Solve constrained allocations}
  {Formulate risky-only and risky/risk-free minimum-variance problems while distinguishing lending, borrowing, and short-sale constraints.}
  {Validate the resulting portfolio}
  {Compare SIM and empirical allocations out of sample without treating a fitted frontier as realized performance.}

\section{From a Return Model to a Portfolio Decision}

The single-index model replaces a dense collection of sample covariances with one market factor and asset-specific residual risks. That parsimony is useful only when it is carried into an optimization problem with compatible units, feasible trading constraints, and honest validation. A smooth fitted frontier is a model output, not evidence that the selected portfolio will perform well.

For $m$ risky assets over one fixed return horizon, let $\boldsymbol\alpha,\boldsymbol\beta\in\R^m$ be SIM intercepts and market loadings, let $\mu_M\in\R$ and $\sigma_M^2\geq0$ be the market return mean and variance, and let $D:=\operatorname{diag}(\sigma_{\varepsilon,1}^2,\ldots,\sigma_{\varepsilon,m}^2)$ contain nonnegative residual variances. The SIM portfolio inputs are
\begin{equation}
  \label{eq:sim-inputs}
  \boldsymbol\mu=\boldsymbol\alpha+\boldsymbol\beta\mu_M,
  \qquad
  \Sigma=\sigma_M^2\boldsymbol\beta\boldsymbol\beta^{\mathsf T}+D.
\end{equation}
Every quantity in \eqnref{sim-inputs} must use the same horizon and return convention. If the notebooks estimate annualized residual variance rates from daily log-growth residuals, the factor variance and expected growth inputs must be annualized consistently.

\begin{proposition}{The SIM covariance is positive semidefinite}{sim-psd}
Let $\boldsymbol\beta\in\R^m$, $\sigma_M^2\geq0$, and let $D\in\R^{m\times m}$ be diagonal with nonnegative entries. Then $\Sigma:=\sigma_M^2\boldsymbol\beta\boldsymbol\beta^{\mathsf T}+D$ is symmetric positive semidefinite because, for every $\mathbf x\in\R^m$,
\begin{equation}
  \label{eq:sim-psd}
  \mathbf x^{\mathsf T}\Sigma\mathbf x
  =\sigma_M^2(\boldsymbol\beta^{\mathsf T}\mathbf x)^2
  +\sum_{i=1}^{m}D_{ii}x_i^2\geq0.
\end{equation}
It is positive definite if every diagonal entry of $D$ is strictly positive.
\end{proposition}

Positive semidefiniteness makes the variance objective convex. It does not establish that the covariance estimate is accurate. Residual cross-correlations, changing betas, and an inadequate market proxy can all produce a valid matrix that is a poor description of future co-movement.

\section{Risky-Only Allocation with SIM Moments}

Let $\mathbf w\in\R^m$ denote risky-asset dollar weights and let $\mu_\star\in\R$ be a target return for the same horizon. A long-only minimum-variance allocation solves
\begin{equation}
  \label{eq:risky-qp}
  \begin{aligned}
  \min_{\mathbf w\in\R^m}\quad &\frac12\mathbf w^{\mathsf T}\Sigma\mathbf w\\
  \text{subject to}\quad
    &\mathbf1^{\mathsf T}\mathbf w=1,\\
    &\boldsymbol\mu^{\mathsf T}\mathbf w\geq\mu_\star,\\
    &\mathbf w\geq\mathbf0.
  \end{aligned}
\end{equation}
The inequality form asks for at least the target. When expected returns contain estimation noise, the optimum may overshoot the target; report the achieved return and solver status rather than assuming equality. Varying $\mu_\star$ traces the modeled constrained boundary. If short sales are permitted, remove or replace the lower bounds and state any leverage or gross-exposure limit.

The rank-one structure also gives a computational shortcut. If $D$ is positive definite, the Sherman--Morrison identity yields
\begin{equation}
  \label{eq:sim-inverse}
  \Sigma^{-1}
  =D^{-1}
  -\frac{\sigma_M^2D^{-1}\boldsymbol\beta\boldsymbol\beta^{\mathsf T}D^{-1}}
  {1+\sigma_M^2\boldsymbol\beta^{\mathsf T}D^{-1}\boldsymbol\beta}.
\end{equation}
Equation~\eqref{eq:sim-inverse} explains why a factor model can scale well, but a constrained quadratic-program solver remains the appropriate tool for \eqnref{risky-qp}.

\notebooklink
  {L8b: SIM Risky-Asset Minimum-Variance Allocation}
  {https://github.com/varnerlab/CHEME-5660-CourseRepository-Fall-2026/blob/main/lectures/week-8/L8b/CHEME-5660-L8b-Example-SIM-MinVar-RA-Fall-2025.ipynb}
  {Construct SIM and empirical covariance inputs for the same ticker universe, solve their constrained frontiers, and compare allocations and realized test-period wealth.}

\section{Adding a Horizon-Matched Risk-Free Asset}

Let $r_f\in\R$ be a known return over the decision horizon, let $w_f\in\R$ be its dollar weight, and let $\mathbf x\in\R^m$ contain risky-asset weights. Because a risk-free payoff has zero covariance with risky returns under the model, only $\mathbf x$ enters portfolio variance. The target-return problem is
\begin{equation}
  \label{eq:risky-riskfree-qp}
  \begin{aligned}
  \min_{\mathbf x,w_f}\quad &\frac12\mathbf x^{\mathsf T}\Sigma\mathbf x\\
  \text{subject to}\quad
    &\mathbf1^{\mathsf T}\mathbf x+w_f=1,\\
    &\boldsymbol\mu^{\mathsf T}\mathbf x+r_fw_f\geq\mu_\star.
  \end{aligned}
\end{equation}
Long-only risky assets impose $\mathbf x\geq\mathbf0$. Risk-free lending imposes $w_f\geq0$; allowing $w_f<0$ represents borrowing at the same rate $r_f$, a strong idealization. Eliminating $w_f=1-\mathbf1^{\mathsf T}\mathbf x$ rewrites the return constraint as
\begin{equation}
  \label{eq:excess-return-constraint}
  (\boldsymbol\mu-r_f\mathbf1)^{\mathsf T}\mathbf x
  \geq\mu_\star-r_f.
\end{equation}
If lending only is allowed, $w_f\geq0$ additionally requires $\mathbf1^{\mathsf T}\mathbf x\leq1$. If borrowing is allowed without limit, that inequality disappears. These are different feasible sets and should never share an unlabeled ``capital allocation line.''

\begin{remark}{When a Treasury STRIP is risk-free}{strip-risk}
A zero-coupon Treasury STRIP delivers a specified nominal payment only when its maturity and the investor's horizon match and the security is held to maturity. Before maturity its market price changes with yields. Inflation changes real purchasing power, and taxes, liquidity, and sovereign-credit assumptions remain. In these notes, ``risk-free'' means a model input with a known nominal horizon payoff, not an asset free of every economic risk.
\end{remark}

\section{The Capital Allocation Geometry and Its Limits}

Let $T$ be a feasible risky portfolio with expected return $\mu_T>r_f$ and standard deviation $\sigma_T>0$. If $y\geq0$ is the total risky weight, combining $y$ in $T$ with $1-y$ in the risk-free asset gives
\begin{equation}
  \label{eq:cal-parameterization}
  \mu_p=r_f+y(\mu_T-r_f),
  \qquad
  \sigma_p=y\sigma_T.
\end{equation}
For unrestricted lending and borrowing at one rate, the efficient risky fund maximizes $(\mu_T-r_f)/\sigma_T$ and \eqnref{cal-parameterization} traces a ray. With lending but no borrowing, only $0\leq y\leq1$ is available on that straight segment; beyond $T$, the constrained efficient set generally returns to the risky-asset boundary. Position limits, different borrowing rates, and short-sale constraints can introduce additional kinks.

\begin{figure}[ht]
  \centering
  \begin{tikzpicture}[x=1.05cm,y=0.78cm,>=Latex,font=\sffamily\small]
    \draw[vnink,line width=0.8pt,-{Latex[length=3mm]}] (0,0)--(8.2,0) node[right] {$\sigma$};
    \draw[vnink,line width=0.8pt,-{Latex[length=3mm]}] (0,0)--(0,5.8) node[above] {$\mu$};
    \draw[vnlink,line width=1.5pt,domain=2.0:6.8,samples=70,smooth]
      plot (\x,{1.65+0.18*(\x-2.0)^2});
    \coordinate (rf) at (0,0.85);
    \coordinate (tan) at (4.65,{1.65+0.18*(4.65-2.0)^2});
    \draw[vncarnelian,line width=1.6pt] (rf)--(tan);
    \draw[vncarnelian,dashed,line width=1.2pt] (tan)--($(rf)!1.55!(tan)$);
    \fill[vnink] (rf) circle (2pt) node[left] {$r_f$};
    \fill[vncarnelian] (tan) circle (2.2pt) node[above left] {$T$};
    \node[text=vncarnelian,rotate=30] at (2.25,2.2) {lending segment};
    \node[text=vncarnelian] at (6.55,4.7) {borrowing extension};
    \node[text=vnlink] at (6.25,2.4) {risky boundary};
  \end{tikzpicture}
  \caption{The solid red segment uses nonnegative risk-free weight: $0\leq y\leq1$. The dashed extension requires borrowing at the modeled risk-free rate. A constrained optimizer must enforce the regime actually available to the investor.}
  \label{fig:constrained-cal}
\end{figure}

The classical separation theorem says investors with common inputs choose the same tangency fund and vary only $y$. This statement relies on mean--variance preferences, common beliefs, one frictionless borrowing/lending rate, and compatible unconstrained opportunities. Once constraints bind, different target returns can select different risky compositions; two-fund separation need not survive globally.

\notebooklink
  {L8b: SIM Risky/Risk-Free Minimum-Variance Allocation}
  {https://github.com/varnerlab/CHEME-5660-CourseRepository-Fall-2026/blob/main/lectures/week-8/L8b/CHEME-5660-L8b-Example-SIM-MinVar-RRFA-Fall-2025.ipynb}
  {Solve risky-only and risky/risk-free frontiers from SIM inputs, identify a candidate tangency allocation, and compare its test-period wealth with market and risk-free benchmarks.}

\notebooklink
  {L8b: Derivation of the SIM Covariance}
  {https://github.com/varnerlab/CHEME-5660-CourseRepository-Fall-2026/blob/main/lectures/week-8/L8b/CHEME-5660-L8b-Derivation-SIM-Covariance-Fall-2025.ipynb}
  {Expand the scalar covariance terms and identify exactly where market-orthogonality and diagonal-residual assumptions enter the result.}

\section{A Defensible Estimation and Validation Workflow}

Use a training window to estimate $\boldsymbol\alpha$, $\boldsymbol\beta$, $D$, $\mu_M$, and $\sigma_M^2$. Freeze those inputs, solve the allocation, and then evaluate the fixed rule on later observations. Re-estimation and rebalancing dates must be specified in advance; otherwise information leakage can make an impossible strategy look successful.

Compare at least three models on the same universe and constraints: the SIM covariance, the full sample covariance, and a simple baseline such as equal weighting or a broad index. Record expected return, modeled variance, realized return and drawdown, turnover, costs, and constraint activity. A lower realized variance in one test window is evidence about that window, not proof of universal dominance.

Parameter uncertainty must also reach the portfolio decision. Refit the SIM over bootstrap or rolling samples, re-solve the optimizer, and inspect the distribution of weights and risk estimates. If small input changes produce radically different positions, constraints, shrinkage, or a simpler portfolio may be more defensible than the nominal optimum.

\keytakeaways
  {In this lecture, we inserted the SIM's low-rank-plus-diagonal moments into constrained risky-only and risky/risk-free allocation problems and separated model geometry from implementable trading rules.}
  {Covariance validity is not forecast validity}
  {The SIM matrix is automatically positive semidefinite, but omitted residual dependence and unstable parameters can still make it a poor future-risk estimate.}
  {Risk-free constraints change the frontier}
  {Lending, borrowing, short sales, and position bounds determine whether the CAL is a segment, a ray, or only part of a kinked efficient boundary.}
  {Validation follows the decision clock}
  {Fit on past data, freeze the portfolio, evaluate later outcomes with costs, and test whether parameter uncertainty destabilizes weights.}
  {Next, allow portfolio weights to evolve through time and distinguish a rebalancing policy from a one-time static optimum.}

\begin{readinglist}
  \item William F. Sharpe, \href{https://doi.org/10.1287/mnsc.9.2.277}{``A Simplified Model for Portfolio Analysis,''} \emph{Management Science} 9(2), 277--293 (1963).
  \item Harry Markowitz, \href{https://www.jstor.org/stable/2975974}{``Portfolio Selection,''} \emph{Journal of Finance} 7(1), 77--91 (1952).
  \item U.S. Department of the Treasury, \href{https://www.treasurydirect.gov/marketable-securities/strips/}{``STRIPS: Separate Trading of Registered Interest and Principal of Securities,''} for the mechanics of stripped Treasury cash flows.
\end{readinglist}

\vfill
{\sffamily\footnotesize\color{vnmuted}\textbf{Educational use and risk notice.} These notes are for instruction only and do not constitute investment advice. Financial decisions involve risk, and model outputs depend on assumptions.}

\end{document}
